\documentclass[10pt,a4paper]{amsart}
\usepackage[margin=19mm]{geometry}
\usepackage{amsmath,amssymb,mathtools}
\usepackage[scr=rsfs,cal=euler]{mathalfa}
\usepackage{xfrac}
\usepackage[unicode,hidelinks,bookmarks=false]{hyperref}
\newcommand{\ie}{i.e.\ }
\newcommand{\cf}{cf.\ }
\newcommand{\resp}{respectively\ }
\newcommand{\Subsection}{\subsection}
\providecommand{\nobreakdash}{\nobreak\hbox{-}}
\setlength{\emergencystretch}{2em}
\multlinegap0pt
\usepackage{bm}
\usepackage{bbm}
\usepackage{dsfont}
\usepackage{xspace}
\usepackage{cancel}
\usepackage{mathtools}
\usepackage{amsthm}
\makeatletter
\@ifundefined{theorem}{\newtheorem{theorem}{Theorem}}{}
\@ifundefined{definition}{\newtheorem{definition}[theorem]{Definition}}{}
\@ifundefined{lemma}{\newtheorem{lemma}[theorem]{Lemma}}{}
\@ifundefined{proposition}{\newtheorem{proposition}[theorem]{Proposition}}{}
\makeatother
\makeatletter
\expandafter\ifx\csname thmx\endcsname\relax
\newenvironment{thmx}{\stepcounter{thm}\begin{thmy}}{\end{thmy}}
\fi
\makeatother
\title[On fast Lyapunov spectra for Markov-R\'{e}nyi maps]{On fast Lyapunov spectra for Markov-R\'{e}nyi maps}
\author{Lulu Fang, Carlos Gustavo Moreira, Zhichao Wang, Yiwei Zhang}
\date{Source: arXiv:2506.16291v1 (CC BY 4.0)}
\begin{document}
\maketitle

\noindent\emph{Source and licence.} On fast Lyapunov spectra for Markov-R\'{e}nyi maps. Version arXiv:2506.16291v1 (CC BY 4.0), distributed under the Creative Commons Attribution 4.0 International licence, as recorded in the arXiv licence metadata for this exact version and shown on the abstract page https://arxiv.org/abs/2506.16291v1. Full rights notice: licenses/arxiv-2506.16291-CC-BY-4.0.txt.\par\smallskip

\noindent\textit{Editorial scope.} Counted unit: the Lemma whose source label is \texttt{l.Hausdorff\_dimension\_Gamma\_infty} in the article (source0.tex lines 1053--1058, source label \texttt{l.Hausdorff\_dimension\_Gamma\_infty}), delivered together with the article's own complete original proof of it (source0.tex lines 1060--1188, one \texttt{proof} environment of 129 lines). No mathematics is written, repaired or re-derived: statement and proof are the article's own text line for line. Required context retained uncounted: def:MR (lines 124--137); p.cylinder (lines 511--519); e.def\_pi\_infty (lines 741--746); l.Hausdorff\_dimension\_limsup\_log\_pi/n (lines 750--755); e.def\_J\_epsilon (lines 774--777). Every definition, display and label of those lines is retained unchanged. Omitted from this package and not redistributed: 1--123, 138--510, 520--740, 747--749, 756--773, 778--1052, 1059--1059, 1189--2527. Nothing inside a retained excerpt is omitted or altered: every retained body is delivered exactly as the article's lines are written, so no omission receipt of any kind accompanies this package. Citations: the retained excerpts carry no citation command at all, so no bibliography entry of the article is transcribed and no reference is redistributed. Reference adaptation: none was needed and none was made; every reference inside the delivered unit resolves inside it, so no unresolved reference or citation is produced. Printed slips kept literal: no printed word, symbol, hypothesis or number of the counted statement or of its proof was altered. Layout and class: the article's own class is \texttt{amsart} with packages including amssymb, latexsym, amsmath, amsxtra, mathrsfs, bm, bookmark, enumerate, hyperref, tikz, geometry, mathtools, graphicx, subcaption; the wrapper is the lane's pinned \texttt{amsart} editorial wrapper at 10pt on A4, which restates the theorem-like environments the retained text needs and adds the article's own macro definitions that the retained text needs (none), with extra wrapper packages (bm, bbm, dsfont, xspace, cancel, mathtools). The delivered numbering is the unit's own. Assets: no retained span contains a figure environment, a graphics-inclusion command, a PDF-page inclusion, a table or a TikZ picture; no image file is copied into this package, the package declares no asset dependency and no image file is redistributed with it. Licence: version arXiv:2506.16291v1 of this article is distributed under the Creative Commons Attribution 4.0 International licence, as recorded in the arXiv licence metadata for this exact version and shown on the abstract page https://arxiv.org/abs/2506.16291v1. A full rights notice is delivered at licenses/arxiv-2506.16291-CC-BY-4.0.txt. Characters: every retained excerpt is pure ASCII, so the delivered engine, XeLaTeX with its default Latin Modern text font, reports no missing character.
\begin{definition}\label{def:MR}
Let $\{I_n\}_{n \in \mathbb N}$ be a collection of disjoint open subintervals of $I=[0,1]$.
A piecewise differentiable interval map $T: \cup_{n \in \mathbb N}\overline{I}_n \to I$ is called a \emph{Markov-R\'{e}nyi map} associated with $\{I_n\}_{n \in \mathbb N}$ if the following conditions hold:
\begin{itemize}
    \item[(1).] for every $n\in \mathbb{N}$, the intervals $I_n$ and $I_{n+1}$ share a common endpoint i.e., $\partial I_{n}\cap\partial I_{n+1}$ is a singleton;
     \item[(2).] for every $n\in \mathbb{N}$, the restriction $T|_{I_n}$ extends to a $C^{1}$ diffeomorphism $T_n$ on an open neighborhood of $\overline{I}_n$;
    \item[(3).] there exists a point $p \in \cup_{n \in \mathbb N}\overline{I}_n$ and an integer $m\in \mathbb{N}$ such that $T(p)=p$, $|T^\prime(p)| \geq 1$, and $|(T^m)^\prime(x)|>1$ for all $x\in \cup_{n \in \mathbb N}\overline{I}_n\setminus \{p\}$;
    \item[(4).] $\overline{T(I_n)}=[0,1]$ for every $n\in \mathbb{N}$;
    \item[(5).] there exist two constants $\gamma>1$ and $C>1$ such that for all $n \in \mathbb{N}$ and $x \in \overline{I}_n$, we have
    \begin{equation*}
        C^{-1}n^{\gamma}\leq |T^\prime_n(x)|\leq Cn^{\gamma}.
    \end{equation*}
\end{itemize}
\end{definition}

\begin{proposition}\label{p.cylinder}
For any $n\in \mathbb N$ and $(i_1,\dots, i_n)\in \mathbb N^n$, the cylinder set $I_n(i_1,\dots,i_n)$ is an interval, and its diameter satisfies
\begin{equation*}\label{e.cylinder}
   \frac{C^{-n}}{(i_1 \cdots i_{n})^{\gamma}}
   \leq |I_n(i_1,\dots,i_n)|
   \leq \frac{C^{n}}{(i_1 \cdots i_{n})^{\gamma}},
\end{equation*}
where the parameters $\gamma$ and $C$ are given in Hypothesis (5) of Definition \ref{def:MR}.
\end{proposition}

        \begin{equation}\label{e.def_pi_infty}
                \Lambda_{\infty}\coloneqq
                \left\{x\in \Lambda:
                \limsup_{n\to\infty}\frac{\sum^n_{k=1}\log a_k(x)}{n}
                =\infty\right\}.
        \end{equation}

        \begin{proposition}\label{l.Hausdorff_dimension_limsup_log_pi/n}
        \begin{equation*}
                \dim_{\mathrm{H}}\Lambda_{\infty}
                =\frac{1}{\gamma}.
        \end{equation*}
\end{proposition}

        \begin{equation}\label{e.def_J_epsilon}
             J_{\epsilon} \coloneqq \sum_{n=1}^{\infty}
        \frac{1}{n^{1+\epsilon}}<\infty.
        \end{equation}

\begin{lemma}\label{l.Hausdorff_dimension_Gamma_infty} For any $1\leq \beta \leq \infty$, we have
        \begin{equation*}\label{e.upperboundGamma_infty}
\dim_{\mathrm{H}}{\Gamma}_{\infty}(\beta)
                \leq \frac{1}{(\gamma-1)\beta+1}.
        \end{equation*}
\end{lemma}

  \begin{proof}
  The proof proceeds by considering three cases: $\beta=1, 1<\beta <\infty$, and $\beta =\infty$.

\textbf{Case I: \bm{$\beta=1.$}} We observe that
\begin{equation*}
            {\Gamma}_{\infty}(1) = \Lambda_{\infty},
        \end{equation*}
        where $\Lambda_{\infty}$ is defined in \eqref{e.def_pi_infty}. It follows from Proposition \ref{l.Hausdorff_dimension_limsup_log_pi/n} that
        \begin{equation}\label{e.dim_Gamma<1/2}
                \dim_{\mathrm{H}}{\Gamma}_{\infty}(1)=\dim_{\mathrm{H}}\Lambda_{\infty}
                =\frac{1}{\gamma},
        \end{equation}
        which completes the proof for this case.


 \textbf{Case II: \bm{$1<\beta<\infty.$}} For any $0<\epsilon<\beta-1$, let $s\coloneqq \frac{1+2\epsilon}{(\gamma-1)(\beta-\epsilon)+1}$. Choose an integer $M\geq 2$ sufficiently large such that
        \begin{equation}\label{e.def_M}
                M^{\epsilon}> 2J_{\epsilon}C^s,
        \end{equation}
        where $J_{\epsilon}$ is defined in \eqref{e.def_J_epsilon}.
For any $n \in \mathbb N$, let
        \begin{equation}\label{e.def_B_n_ep_M}
                B_n(\epsilon,M)\coloneqq
                \Big\{
                x\in  \Lambda:
                a_{n+1}(x)>\big(
                        \prod_{k=1}^{n}
                        a_k(x)
                        \big)^{\beta-1-\epsilon}\ \text{and}\
                \prod_{k=1}^{n}
                a_k(x)\geq M^n
                \Big\}.
                \end{equation}
   Thus,
        \begin{equation}\label{e.cover_a_n}
                {\Gamma}_{\infty}(\beta)
                \subset \bigcap_{N=1}^{\infty}
                \bigcup_{n=N}^{\infty}
                B_n(\epsilon,M).
        \end{equation}
For any $n \in \mathbb N$, define
        \begin{equation}\label{e.def_D_n}
            \mathcal{D} _n(M)\coloneqq
        \left\{(i_1,\dots, i_{n})\in
        \mathbb{N}^n:
        i_1\cdots i_{n}\geq M^n\right\},
        \end{equation}
 and for any $(i_1,\dots, i_{n})
        \in \mathcal{D} _n(M)$, let
        \begin{equation}\label{e.def_J_n_another}
                J_n(i_1,\dots,i_{n}) \coloneqq
                \bigcup_{k>(i_1\cdots i_{n})
                ^{\beta-1-\epsilon}} I_{n+1}(i_1,\dots,i_{n},k).
        \end{equation}
Then $B_n(\epsilon,M)$ can be reformulated as
        \begin{equation}\label{e.a_n_J_n}
                B_n(\epsilon,M)=
                \bigcup_{(i_1,\dots, i_{n})
                \in \mathcal{D} _n(M)}
                J_n(i_1,\dots,i_{n}).
        \end{equation}
 Combining this with \eqref{e.cover_a_n}, for every $N\geq 1$, we see that,
        \begin{equation}\label{e.cover_by_J}
                {\Gamma}_{\infty}(\beta)\subset
                \bigcup^\infty_{n=N}\bigcup_{(i_0,\dots, i_{n-1})
                \in \mathcal{D}_n(M)}J_n(i_0,\dots,i_{n-1}).
        \end{equation}
For any $n \in \mathbb N$ and $(i_1,\dots, i_{n})\in \mathcal{D} _n(M)$, we estimate the diameter of $J_n(i_1,\dots,i_{n})$. Since
$$
\sum_{k>(i_1\cdots i_{n})
        ^{\beta-1-\epsilon}}\frac{1}{k^{\gamma}} \leq \frac{2^{\gamma-1}}
        {(\gamma-1)(i_1\cdots i_{n})^{(\beta-1-\epsilon)(\gamma-1)}},
$$
we deduce from Proposition \ref{p.cylinder} and \eqref{e.def_J_n_another} that
     \begin{alignat*}{2}
        |J_n(i_1,\dots,i_{n})|\leq\frac{C^{n+1}}
        {(i_1\cdots i_{n})^{\gamma}}
        \sum_{k>(i_1\cdots i_{n})
        ^{\beta-1-\epsilon}}\frac{1}{k^{\gamma}}
       \leq \frac{C^{n+1}2^{\gamma-1}}
        {(\gamma-1)(i_1\cdots i_{n})^{(\gamma-1)(\beta-\epsilon)+1}}.
        \end{alignat*}
Note that $i_1\cdots i_{n}\geq M^n$, we derive that
 \begin{equation*}\label{J^s}
|J_n(i_1,\dots,i_{n})|^s \leq \left(\frac{2^{\gamma-1}C}{\gamma-1}\right)^s \frac{C^{sn}}
        {(i_1\cdots i_{n})^{1+2\epsilon}} \leq \left(\frac{2^{\gamma-1}C}{\gamma-1}\right)^s \left(\frac{C^{s}}{M^\epsilon}\right)^n \frac{1}
        {(i_1\cdots i_{n})^{1+\epsilon}}.
 \end{equation*}
 Hence,
\begin{alignat*}{2}
 \sum_{(i_1,\dots, i_{n})\in \mathcal{D} _n(M)}|J_n(i_1,\dots,i_{n})|^s \leq &\left(\frac{2^{\gamma-1}C}{\gamma-1}\right)^s \left(\frac{C^{s}}{M^\epsilon}\right)^n \sum_{(i_1,\dots, i_{n})\in \mathcal{D} _n(M)} \frac{1}
        {(i_1\cdots i_{n})^{1+\epsilon}}\\
\leq & \left(\frac{2^{\gamma-1}C}{\gamma-1}\right)^s \left(\frac{J_\epsilon C^{s}}{M^\epsilon}\right)^n.
\end{alignat*}
From the definition of the $s$-dimensional Hausdorff measure, it follows that
 \begin{alignat*}{2}
 \mathcal{H}^{s}\left({\Gamma}_{\infty}(\beta)\right)
                \leq& \liminf_{N \to \infty}
                \sum_{n=N}^{\infty}
                \sum_{(i_1,\dots, i_{n})
                \in \mathcal{D} _n(M)}
                |J_n(i_1,\dots,i_{n})|^s \\
                \leq & \left(\frac{2^{\gamma-1}C}{\gamma-1}\right)^s \liminf_{N \to \infty}
                \sum_{n=N}^{\infty}\left(\frac{J_\epsilon C^{s}}{M^\epsilon}\right)^n\\
                 \leq & \left(\frac{2^{\gamma-1}C}{\gamma-1}\right)^s \liminf_{N \to \infty}
                \sum_{n=N}^{\infty} \frac{1}{2^n} =0. && \  \big(\text{by}\ \eqref{e.def_M}\big)
        \end{alignat*}
This implies that
 \begin{equation*}
                \dim_{\mathrm{H}}
              {\Gamma}_{\infty}(\beta)
                \leq s\coloneqq \frac{1+2\epsilon}{(\gamma-1)(\beta-\epsilon)+1}.
        \end{equation*}
 Letting $\epsilon\to 0^{+}$, we obtain $$\dim_{\mathrm{H}}{\Gamma}_{\infty}
        (\beta)\leq \frac{1}{(\gamma-1)\beta+1}.$$


\textbf{Case III: \bm{$\beta=\infty.$}} For any $K>1$, since
  \begin{equation*}
           {\Gamma}_{\infty}(\infty)\subset
           {\Gamma}_{\infty}(K),
        \end{equation*}
it follows from the result of Case II that
        \begin{equation*}
            \dim_{\mathrm{H}}{\Gamma}_{\infty}(\infty)
            \leq \dim_{\mathrm{H}}{\Gamma}_{\infty}(K)=\frac{1}{(\gamma-1)K+1}.
        \end{equation*}
Taking the limit as $K\to \infty$, we have $\dim_{\mathrm{H}}{\Gamma}_{\infty}(\infty)=0$.
\end{proof}

\end{document}
